\section{Introduction}

\subsection{Cancer as a nonequilibrium system}
Tumor growth is a driven, open system with matter and energy fluxes. Nonequilibrium physics provides tools to interpret multistability, thresholds, and pattern formation in tumor--immune competition \cite{murray2002}.

\subsection{Limitations of reciprocal population models}
Symmetric Lotka--Volterra or logistic models miss thresholds, hysteresis, and asymmetric suppression/activation between cancer, healthy tissue, and immune cells. These limitations motivate adding Allee effects and nonreciprocal couplings.

\subsection{Allee effects and critical thresholds}
An Allee term introduces a viability threshold for cancer: below it, extinction; above it, proliferation. Weak Allee yields a soft slowdown; strong Allee creates a sharp barrier. Autocrine signalling can induce Allee-like behaviour in tumors \cite{gerlee2022}; phenotypic switching can reinforce extinction for subcritical populations \cite{bottger2015}.

\subsection{Nonreciprocal interactions in cancer--immune dynamics}
Suppression of immunity by tumor cells and cytotoxic action of immune cells are not symmetric. Including nonreciprocity alters nullclines, fixed points, and transient routes \cite{delitala2007,gatenby2004}. We treat nonreciprocity as asymmetry in interaction coefficients and as an external control field.

\subsection{Summary of contributions}
\begin{itemize}
    \item Formulate ODE/PDE models with weak/strong Allee effects, nonreciprocity, and an immunological control field $u(x,t)$.
    \item Characterize fixed points and linear stability; identify Allee-driven bifurcations in mean field.
    \item Extend to reaction--diffusion dynamics; report spatial patterns, correlation lengths, and spectral content.
    \item Evaluate control strategies (binary vs adaptive) and their cost--efficiency trade-offs across weak/strong Allee regimes.
    \item Provide figure set: Fig.~1 conceptual scheme; Fig.~2 phase planes/nullclines; Fig.~3 bifurcations; Fig.~4--5 spatial patterns and spectra; Fig.~6 control field; Table~I parameters; Table~II cost vs suppression.
\end{itemize}

% Figure 1 placeholder: conceptual scheme of cancer--immune interactions and control.
\begin{figure}[htbp]
    \centering
    \includegraphics[draft=false,width=0.55\linewidth]{steady.png}
    \caption{Conceptual phase portrait (weak Allee, $\mu=1$) illustrating nullclines and an unstable fixed point as threshold between extinction and proliferation.}
\end{figure}

